\documentclass[12pt,a4paper]{article}
\usepackage{fontspec}
\usepackage{polyglossia}
\setmainfont{TeX Gyre Termes}
\newfontfamily\cyrillicfont{TeX Gyre Termes}
\setdefaultlanguage{ukrainian}
\setotherlanguage{english}
\usepackage{amsmath,amssymb,geometry,graphicx}
\geometry{margin=25mm}
\setlength{\parindent}{1.25cm}
\setlength{\parskip}{0.25em}
\begin{document}
\title{2022 Interpolation of different types of spiral-like curves by gaus-interpolation methods (ст ВАК)}
\maketitle
2022 Interpolation of different types of spiral-like curves by gaus-interpolation methods (￿￿ ￿￿￿) 27 ￿￿￿￿￿￿ 2026 ￿. 2022 Interpolation of different types of spiral-like curves by gaus-interpolation methods (￿￿ ￿￿￿) 27 ￿￿￿￿￿￿ 2026 ￿. ORCID: https://orcid.org/0000-0002-1953-0410, Fundamental Problems in Computer Science ISSN 2706-8145, Control systems and computers, 2022, № 3 3 DOI https://doi.org/10.15407/csc.2022.03.003 UDC 514.18 I U .V . SYDORENKO, PhD (Eng. Sc.), Assistant Professor, National Technical University of Ukraine “Igor Sikorsky Kyiv Polytechnic Institute”, 37, Prosp. Peremohy, Kyiv, Ukraine, 03056, suliko3@ukr.net A.I. ONYSKO, Ph.D. (Mil. Sc.), Associate Professor, National Technical University of Ukraine “Igor Sikorsky Kyiv Polytechnic Institute”, 37, Prosp. Peremohy, Kyiv, Ukraine, 03056, ORCID: https://orcid.org/0000- 0002-1953-0410, kw\_fedun@ukr.net O.V . SHALDENKO, PhD (Eng. Sc.), National Technical University of Ukraine “Igor Sikorsky Kyiv Polytechnic Institute”, 37, Prosp. Peremohy, Kyiv, Ukraine, 03056, ORCID: https://orcid.org/0000- 0001-6730-965X, o.shaldenko@gmail.com M.V . HORODETSKYI, PhD Student, National Technical University of Ukraine “Igor Sikorsky Kyiv Polytechnic Institute”, 37, Prosp. Peremohy, Kyiv, Ukraine, 03056, ORCID: https://orcid.org/0000- 0003-4673-3894, horodetkyimykola@gmail.com INTERPOLATION OF DIFFERENT TYPES OF SPIRAL-LIKE CURVES BY GAUS-INTERPOLATION METHODS The methods of Gaussian interpolation of helical curves are studied in the article. A comparative analysis with the standard Lagrange method was carried out. The results of the work of Gaussian methods with different types of interpolation steps are demonstrated. Recommendations on the optimal choice of method for spiral curves are offered. An analysis of each of the Gauss interpolation methods on the most common types of spirals has been carried out. A solution for input data with a non-constant step is proposed. Keywords: interpolation, Gaussian interpolation function, spiral curves, interpolation step, interpolation error. Introduction Nowadays, more and more systems for proces- sing large data arrays based on modern information technologies are being created. The need to deve- lop new effective methods for analysing, display- ing, storing, and processing empirical data arise with an increase in the amount of information and the need to process it. Interpolation of spiral curves causes scientific and practical enthusiasm in the design of various technical means. In fact there is different spheres of life there are such curves: in the surrounding nature, in agricultural machinery , in art, in space- craft, in military equipment, etc. The spiral ap- pears and is studied, and then these results are used in the field of aircraft and shipbuilding, construc- Iu.V . Sydorenko, A.I. Onysko, O.V . Shaldenko, M.V . Horodetskyi 4 ISSN 2706-8145, ￿￿￿￿￿￿￿ ￿￿￿￿￿￿￿￿￿ ￿￿ ￿￿￿￿'￿￿￿￿￿, 2022, № 3 tion, information encryption, designing railways and roads, encoding confidential information, cal- culating the movement of robots and drones, and more. Therefore, the problem of interpolation of spiral curves while minimizing the calculation er- ror arose. When creating a system for interpolating spiral curves, the advantages and disadvantages of classi- cal interpolation methods were taken into account. The article provides a comparative analysis of 1
\par
the Gaussian interpolation of spiral curves with inter- polation by the Lagrange 1 2
\par
\end{document}
